% Topic 17: Division, Span and Style Attribute
\section{Division, Span ও Style Attribute}

\begin{learningbox}
এই টপিক শেষে শিক্ষার্থী পারবে:
\begin{itemize}
\item \texttt{<div>} tag-এর কাজ ও ব্যবহার ব্যাখ্যা করতে।
\item \texttt{<span>} tag-এর কাজ ও ব্যবহার ব্যাখ্যা করতে।
\item \texttt{<div>} ও \texttt{<span>}-এর পার্থক্য লিখতে।
\item \texttt{style} attribute ব্যবহার করে inline CSS লিখতে।
\item \texttt{color}, \texttt{background-color}, \texttt{font-size}, \texttt{text-align}, \texttt{margin}, \texttt{padding}, \texttt{border} ইত্যাদি basic CSS property ব্যবহার করতে।
\end{itemize}
\end{learningbox}

আগের topic-গুলোতে আমরা HTML দিয়ে heading, paragraph, list, hyperlink, image ও table তৈরি করেছি। এখন ওয়েবপেজের বিভিন্ন অংশকে group করা এবং নির্দিষ্ট অংশে style প্রয়োগ করা শিখব। এই topic-এ তিনটি বিষয় সবচেয়ে গুরুত্বপূর্ণ: \texttt{<div>}, \texttt{<span>} এবং \texttt{style} attribute।

\subsection{\texttt{<div>} Tag কী?}

\begin{definitionbox}{\texttt{<div>} tag}
HTML document-এর বিভিন্ন অংশ বা section-কে আলাদা block আকারে group করার জন্য ব্যবহৃত container tag-কে \texttt{<div>} tag বলে।
\end{definitionbox}

\texttt{<div>} হলো HTML-এর একটি generic container tag। এটি webpage-এর বড় অংশ, section, layout block বা একাধিক element একসাথে group করার জন্য ব্যবহার করা হয়। নিজে নিজে এটি বিশেষ design তৈরি করে না; তবে CSS style দিলে এটি section, box বা layout তৈরিতে খুব গুরুত্বপূর্ণ হয়ে ওঠে।

\subsubsection{Syntax}

\begin{examplebox}{Code}
\begin{verbatim}
<div>
  Content here
</div>
\end{verbatim}
\end{examplebox}

\subsubsection{Example}

\begin{examplebox}{Code}
\begin{verbatim}
<div>
  <h2>About College</h2>
  <p>ABC College is a reputed institution.</p>
</div>
\end{verbatim}
\end{examplebox}

এখানে heading ও paragraph একসাথে একটি section বা block হিসেবে group করা হয়েছে।

\subsection{\texttt{<div>} Tag-এর ব্যবহার}

\texttt{<div>} সাধারণত webpage-এর বড় অংশ ভাগ করতে ব্যবহৃত হয়।

\begin{itemize}
\item Header section তৈরি করতে
\item Navigation/menu area group করতে
\item Main content section তৈরি করতে
\item Sidebar তৈরি করতে
\item Footer section তৈরি করতে
\item একাধিক element একসাথে style করতে
\item webpage layout তৈরি করতে
\end{itemize}

\subsection{Example: Page Section}

\begin{examplebox}{Code}
\begin{verbatim}
<!DOCTYPE html>
<html>
<head>
  <title>Div Example</title>
</head>
<body>
  <div style="background-color:lightblue; padding:10px;">
    <h1>ABC College</h1>
    <p>Welcome to our official website.</p>
  </div>

  <div style="background-color:lightyellow; padding:10px;">
    <h2>Notice</h2>
    <p>HSC ICT class will start from Sunday.</p>
  </div>
</body>
</html>
\end{verbatim}
\end{examplebox}

\begin{examplebox}{Output ধারণা}
প্রথম \texttt{<div>} একটি light blue background section তৈরি করবে। দ্বিতীয় \texttt{<div>} একটি light yellow background section তৈরি করবে।
\end{examplebox}

\subsection{\texttt{<span>} Tag কী?}

\begin{definitionbox}{\texttt{<span>} tag}
HTML document-এর text বা inline content-এর নির্দিষ্ট অংশকে আলাদা করে style বা script প্রয়োগ করার জন্য ব্যবহৃত inline container tag-কে \texttt{<span>} tag বলে।
\end{definitionbox}

\texttt{<span>} হলো HTML-এর একটি inline container tag। এটি কোনো paragraph বা line-এর ভিতরে text-এর নির্দিষ্ট অংশ আলাদা করে style করার জন্য ব্যবহৃত হয়। \texttt{<span>} সাধারণত নতুন line তৈরি করে না।

\subsubsection{Syntax}

\begin{examplebox}{Code}
\begin{verbatim}
<span>Text</span>
\end{verbatim}
\end{examplebox}

\subsubsection{Example}

\begin{examplebox}{Code}
\begin{verbatim}
<p>This is normal text.
  <span style="color:red;">This is red text.</span>
</p>
\end{verbatim}
\end{examplebox}

এখানে শুধু \textbf{This is red text.} অংশটি red color হবে।

\subsection{\texttt{<div>} ও \texttt{<span>}-এর পার্থক্য}

\begin{center}
\begin{tabular}{|p{0.24\textwidth}|p{0.30\textwidth}|p{0.30\textwidth}|}
\hline
\textbf{বিষয়} & \textbf{\texttt{<div>}} & \textbf{\texttt{<span>}} \\
\hline
Element type & Block-level container & Inline-level container \\
\hline
কাজ & বড় section বা block group করে & text/inline অংশ group করে \\
\hline
নতুন line তৈরি করে? & সাধারণত হ্যাঁ & না \\
\hline
ব্যবহার & layout, section, card, box & paragraph-এর ভিতরে নির্দিষ্ট word/phrase style \\
\hline
Example & header, footer, notice box & red word, highlighted text \\
\hline
Closing tag & আছে & আছে \\
\hline
\end{tabular}
\end{center}

\begin{keypointbox}{সহজ মনে রাখার নিয়ম}
\texttt{div} = division = বড় ভাগ বা section। \texttt{span} = ছোট inline অংশ।
\end{keypointbox}

\subsection{Style Attribute কী?}

\begin{definitionbox}{style attribute}
HTML element-এর opening tag-এর মধ্যে সরাসরি CSS property ও value ব্যবহার করে style প্রয়োগ করার পদ্ধতিকে style attribute বলা হয়।
\end{definitionbox}

\texttt{style} হলো HTML-এর একটি global attribute। এর মাধ্যমে কোনো element-এ সরাসরি CSS style দেওয়া যায়।

\subsubsection{Basic Syntax}

\begin{examplebox}{Code}
\begin{verbatim}
<tag style="property:value;">Content</tag>
\end{verbatim}
\end{examplebox}

\subsubsection{Example}

\begin{examplebox}{Code}
\begin{verbatim}
<p style="color:red;">This text is red.</p>
\end{verbatim}
\end{examplebox}

\begin{center}
\begin{tabular}{|p{0.25\textwidth}|p{0.50\textwidth}|}
\hline
\textbf{অংশ} & \textbf{পরিচয়} \\
\hline
\texttt{style} & attribute \\
\hline
\texttt{color} & CSS property \\
\hline
\texttt{red} & property value \\
\hline
\texttt{color:red;} & CSS declaration \\
\hline
\end{tabular}
\end{center}

\subsection{Style Attribute-এর সঠিক নিয়ম}

\texttt{style} attribute-এর ভিতরে CSS property ও value লিখতে হয়।

\subsubsection{সঠিক structure}

\begin{verbatim}
style="property:value; property:value;"
\end{verbatim}

\subsubsection{Example}

\begin{examplebox}{Code}
\begin{verbatim}
<div style="color:red; background-color:yellow;">
  HSC ICT
</div>
\end{verbatim}
\end{examplebox}

\begin{boardbox}{গুরুত্বপূর্ণ নিয়ম}
\begin{itemize}
\item Property ও value-এর মাঝে colon \texttt{:} দিতে হয়।
\item একাধিক declaration-এর মাঝে semicolon \texttt{;} দিতে হয়।
\item Attribute value quotation mark-এর মধ্যে লেখা ভালো।
\item CSS property-value লেখার সময় \texttt{=} ব্যবহার করা যাবে না।
\item comma দিয়ে CSS declaration আলাদা করা যাবে না।
\end{itemize}
\end{boardbox}

\subsection{Correct vs Incorrect Style Syntax}

\subsubsection{Correct}

\begin{examplebox}{Code}
\begin{verbatim}
<div style="color:red; float:left;">
  ICT
</div>
\end{verbatim}
\end{examplebox}

\subsubsection{Incorrect}

\begin{examplebox}{Code}
\begin{verbatim}
<div style="color=red, float=left">
  ICT
</div>
\end{verbatim}
\end{examplebox}

এটি ভুল, কারণ \texttt{color=red} নয়, সঠিক হলো \texttt{color:red}। আবার \texttt{float=left} নয়, সঠিক হলো \texttt{float:left}। CSS declaration আলাদা করতে comma নয়, semicolon ব্যবহার করতে হয়।

\subsection{Inline CSS কী?}

\begin{definitionbox}{Inline CSS}
কোনো HTML element-এর opening tag-এর মধ্যে \texttt{style} attribute ব্যবহার করে সরাসরি CSS property-value লেখা হলে তাকে Inline CSS বলে।
\end{definitionbox}

\begin{examplebox}{Code}
\begin{verbatim}
<h1 style="color:blue;">HSC ICT</h1>
<p style="font-size:20px;">Web Design and HTML</p>
\end{verbatim}
\end{examplebox}

\subsection{গুরুত্বপূর্ণ CSS Property}

HSC level-এ style attribute শেখার জন্য নিচের propertyগুলো বেশি গুরুত্বপূর্ণ।

\begin{center}
\begin{tabular}{|p{0.26\textwidth}|p{0.34\textwidth}|p{0.26\textwidth}|}
\hline
\textbf{Property} & \textbf{কাজ} & \textbf{Example} \\
\hline
\texttt{color} & text color নির্ধারণ করে & \texttt{color:red;} \\
\hline
\texttt{background-color} & background color নির্ধারণ করে & \texttt{background-color:yellow;} \\
\hline
\texttt{font-size} & font size নির্ধারণ করে & \texttt{font-size:20px;} \\
\hline
\texttt{font-family} & font style/family নির্ধারণ করে & \texttt{font-family:Arial;} \\
\hline
\texttt{text-align} & text alignment নির্ধারণ করে & \texttt{text-align:center;} \\
\hline
\texttt{width} & element-এর প্রস্থ নির্ধারণ করে & \texttt{width:200px;} \\
\hline
\texttt{height} & element-এর উচ্চতা নির্ধারণ করে & \texttt{height:100px;} \\
\hline
\texttt{margin} & element-এর বাইরের ফাঁকা জায়গা & \texttt{margin:10px;} \\
\hline
\texttt{padding} & element-এর ভিতরের ফাঁকা জায়গা & \texttt{padding:10px;} \\
\hline
\texttt{border} & border তৈরি করে & \texttt{border:1px solid black;} \\
\hline
\texttt{float} & element-কে left/right দিকে ভাসায় & \texttt{float:left;} \\
\hline
\end{tabular}
\end{center}

\subsection{\texttt{<div>}-এ Style Attribute ব্যবহার}

\begin{examplebox}{Code}
\begin{verbatim}
<!DOCTYPE html>
<html>
<head>
  <title>Div Style Example</title>
</head>
<body>
  <div style="background-color:lightyellow; border:1px solid black; padding:10px; width:300px;">
    <h2>Notice</h2>
    <p>ICT class will start at 10:00 AM.</p>
  </div>
</body>
</html>
\end{verbatim}
\end{examplebox}

\begin{examplebox}{Output ধারণা}
একটি notice box দেখা যাবে। এর background light yellow হবে, border থাকবে, ভিতরে padding থাকবে এবং box-এর ভিতরে Notice heading ও paragraph থাকবে।
\end{examplebox}

\subsection{\texttt{<span>}-এ Style Attribute ব্যবহার}

\begin{examplebox}{Code}
\begin{verbatim}
<!DOCTYPE html>
<html>
<head>
  <title>Span Style Example</title>
</head>
<body>
  <p>
    HTML is used to create
    <span style="color:red; font-weight:bold;">web pages</span>.
  </p>
</body>
</html>
\end{verbatim}
\end{examplebox}

\begin{examplebox}{Output ধারণা}
পুরো paragraph normal থাকবে, কিন্তু \textbf{web pages} অংশটি red ও bold হবে।
\end{examplebox}

\subsection{\texttt{<div>} ও \texttt{<span>} একসাথে ব্যবহার}

\begin{examplebox}{Code}
\begin{verbatim}
<!DOCTYPE html>
<html>
<head>
  <title>Div and Span Example</title>
</head>
<body>
  <div style="background-color:#eeeeee; padding:15px; border:1px solid gray;">
    <h1>HSC ICT Book</h1>
    <p>
      This chapter is about
      <span style="color:blue; font-weight:bold;">Web Design and HTML</span>.
    </p>
  </div>
</body>
</html>
\end{verbatim}
\end{examplebox}

\begin{keypointbox}{Explanation}
\texttt{<div>} পুরো section-টিকে gray background box করেছে। \texttt{<span>} শুধু ``Web Design and HTML'' অংশে blue bold style দিয়েছে।
\end{keypointbox}

\subsection{Margin ও Padding}

\subsubsection{Margin কী?}

\texttt{margin} element-এর বাইরের ফাঁকা জায়গা নির্ধারণ করে।

\begin{examplebox}{Code}
\begin{verbatim}
<div style="margin:20px;">Content</div>
\end{verbatim}
\end{examplebox}

\subsubsection{Padding কী?}

\texttt{padding} element-এর content ও border-এর মাঝের ভিতরের ফাঁকা জায়গা নির্ধারণ করে।

\begin{examplebox}{Code}
\begin{verbatim}
<div style="padding:20px;">Content</div>
\end{verbatim}
\end{examplebox}

\begin{center}
\begin{tabular}{|p{0.22\textwidth}|p{0.31\textwidth}|p{0.31\textwidth}|}
\hline
\textbf{বিষয়} & \textbf{Margin} & \textbf{Padding} \\
\hline
অবস্থান & border-এর বাইরে & border-এর ভিতরে \\
\hline
কাজ & element-এর বাইরে space তৈরি করে & content-এর চারপাশে ভিতরের space তৈরি করে \\
\hline
Example & \texttt{margin:10px;} & \texttt{padding:10px;} \\
\hline
\end{tabular}
\end{center}

\subsection{Border}

Element-এর চারপাশে border দিতে \texttt{border} property ব্যবহার করা হয়।

\begin{verbatim}
style="border:1px solid black;"
\end{verbatim}

\begin{examplebox}{Code}
\begin{verbatim}
<div style="border:2px solid blue; padding:10px;">
  Important Note
</div>
\end{verbatim}
\end{examplebox}

\begin{center}
\begin{tabular}{|p{0.20\textwidth}|p{0.48\textwidth}|}
\hline
\textbf{অংশ} & \textbf{অর্থ} \\
\hline
\texttt{2px} & border-এর পুরুত্ব \\
\hline
\texttt{solid} & border-এর ধরন \\
\hline
\texttt{blue} & border color \\
\hline
\end{tabular}
\end{center}

\subsection{Text Align}

Text বাম, মাঝখান বা ডানে রাখতে \texttt{text-align} property ব্যবহার করা হয়।

\begin{examplebox}{Code}
\begin{verbatim}
<h1 style="text-align:center;">ABC College</h1>
\end{verbatim}
\end{examplebox}

\begin{center}
\begin{tabular}{|p{0.25\textwidth}|p{0.50\textwidth}|}
\hline
\textbf{Value} & \textbf{কাজ} \\
\hline
\texttt{left} & বামে রাখে \\
\hline
\texttt{center} & মাঝখানে রাখে \\
\hline
\texttt{right} & ডানে রাখে \\
\hline
\texttt{justify} & দুই পাশে সমানভাবে ছড়িয়ে দেয় \\
\hline
\end{tabular}
\end{center}

\subsection{Color Value}

HTML/CSS-এ color কয়েকভাবে দেওয়া যায়।

\begin{center}
\begin{tabular}{|p{0.25\textwidth}|p{0.50\textwidth}|}
\hline
\textbf{পদ্ধতি} & \textbf{Example} \\
\hline
Color name & \texttt{red}, \texttt{blue}, \texttt{green} \\
\hline
Hex code & \texttt{\#FF0000}, \texttt{\#FFFFFF} \\
\hline
RGB & \texttt{rgb(255, 0, 0)} \\
\hline
\end{tabular}
\end{center}

\begin{boardbox}{MCQ point}
\texttt{\#FFFFFF} সাদা color নির্দেশ করে।
\end{boardbox}

\subsection{Internal CSS: \texttt{<style>} Tag}

যদি অনেক element-এ একই style দিতে হয়, তাহলে বারবার inline style না লিখে \texttt{<head>} section-এর ভিতরে \texttt{<style>} tag ব্যবহার করা যায়।

\begin{examplebox}{Code}
\begin{verbatim}
<!DOCTYPE html>
<html>
<head>
  <title>Internal CSS Example</title>
  <style>
    div {
      background-color: lightblue;
      padding: 10px;
      margin: 10px;
      border: 1px solid black;
    }
  </style>
</head>
<body>
  <div>First Box</div>
  <div>Second Box</div>
</body>
</html>
\end{verbatim}
\end{examplebox}

\begin{examplebox}{Output ধারণা}
দুইটি \texttt{<div>} একই style পাবে।
\end{examplebox}

\subsection{Inline CSS বনাম Internal CSS}

\begin{center}
\begin{tabular}{|p{0.21\textwidth}|p{0.31\textwidth}|p{0.31\textwidth}|}
\hline
\textbf{বিষয়} & \textbf{Inline CSS} & \textbf{Internal CSS} \\
\hline
কোথায় লেখা হয় & element-এর opening tag-এ & \texttt{<head>}-এর ভিতরে \texttt{<style>} tag-এ \\
\hline
Syntax & \texttt{<p style="color:red;">} & \texttt{p \{ color:red; \}} \\
\hline
ব্যবহার & একক element দ্রুত style করা & একটি page-এর অনেক element style করা \\
\hline
সুবিধা & সহজ ও দ্রুত & code পরিষ্কার থাকে \\
\hline
সীমাবদ্ধতা & code messy হয় & CSS syntax জানা দরকার \\
\hline
\end{tabular}
\end{center}

\subsection{Full Example: Mini Webpage Layout}

\begin{examplebox}{Code}
\begin{verbatim}
<!DOCTYPE html>
<html>
<head>
  <title>Mini Webpage Layout</title>
</head>
<body>
  <div style="background-color:#004080; color:white; padding:15px; text-align:center;">
    <h1>ABC College</h1>
    <p>Official Website</p>
  </div>

  <div style="background-color:#eeeeee; padding:10px; margin-top:10px;">
    <a href="index.html">Home</a> |
    <a href="about.html">About</a> |
    <a href="notice.html">Notice</a> |
    <a href="contact.html">Contact</a>
  </div>

  <div style="padding:15px; border:1px solid #cccccc; margin-top:10px;">
    <h2>Welcome</h2>
    <p>
      This website is designed for
      <span style="color:red; font-weight:bold;">HSC ICT students</span>.
    </p>
  </div>

  <div style="background-color:#222222; color:white; padding:10px; text-align:center; margin-top:10px;">
    <p>&copy; 2026 ABC College</p>
  </div>
</body>
</html>
\end{verbatim}
\end{examplebox}

\begin{examplebox}{Output Explanation}
প্রথম \texttt{<div>} header section, দ্বিতীয় \texttt{<div>} navigation section, তৃতীয় \texttt{<div>} main content section এবং শেষ \texttt{<div>} footer section। \texttt{<span>} দিয়ে নির্দিষ্ট text red bold করা হয়েছে।
\end{examplebox}

\subsection{HSC Board/Exam Focus}

\begin{boardbox}{বেশি গুরুত্বপূর্ণ}
\begin{itemize}
\item Document-এর বিভিন্ন অংশ আলাদা করতে \texttt{<div>} ব্যবহার হয়।
\item নির্দিষ্ট inline অংশ নির্বাচন করতে \texttt{<span>} ব্যবহার হয়।
\item \texttt{<div>} block-level, \texttt{<span>} inline-level।
\item \texttt{style} attribute দিয়ে inline CSS লেখা হয়।
\item Style syntax: \texttt{property:value;}
\item Multiple style declaration semicolon \texttt{;} দিয়ে আলাদা হয়।
\item \texttt{color:red; float:left;} সঠিক style syntax।
\item \texttt{color=red} ভুল syntax।
\item \texttt{background-color}, \texttt{color}, \texttt{width}, \texttt{height}, \texttt{margin}, \texttt{padding}, \texttt{float} গুরুত্বপূর্ণ property।
\end{itemize}
\end{boardbox}

\subsection{Exam-style Short Answer}

\begin{enumerate}
\item \textbf{\texttt{<div>} tag কী?}\\
উত্তর: HTML document-এর বিভিন্ন অংশ বা section-কে আলাদা block আকারে group করার জন্য ব্যবহৃত container tag-কে \texttt{<div>} tag বলে। এটি header, menu, content, sidebar বা footer section তৈরি করতে ব্যবহার করা যায়।

\item \textbf{\texttt{<span>} tag কী?}\\
উত্তর: HTML document-এর text বা inline content-এর নির্দিষ্ট অংশকে আলাদা করে style বা script প্রয়োগ করার জন্য ব্যবহৃত inline container tag-কে \texttt{<span>} tag বলে। এটি paragraph-এর ভিতরে কোনো নির্দিষ্ট word বা phrase style করতে ব্যবহৃত হয়।

\item \textbf{\texttt{<div>} ও \texttt{<span>}-এর পার্থক্য লেখ।}\\
উত্তর: \texttt{<div>} একটি block-level container, যা webpage-এর বড় section বা block group করতে ব্যবহৃত হয়। অন্যদিকে \texttt{<span>} একটি inline-level container, যা text বা line-এর ভিতরে নির্দিষ্ট অংশ group করতে ব্যবহৃত হয়।

\item \textbf{Style attribute কী?}\\
উত্তর: HTML element-এর opening tag-এর মধ্যে সরাসরি CSS property ও value ব্যবহার করে style প্রয়োগ করার পদ্ধতিকে style attribute বলে। যেমন \texttt{<p style="color:red;">ICT</p>}।

\item \textbf{\texttt{style="color:red; float:left;"} কেন সঠিক?}\\
উত্তর: কারণ style attribute-এর ভিতরে CSS property ও value colon \texttt{:} দিয়ে লেখা হয় এবং একাধিক declaration semicolon \texttt{;} দিয়ে আলাদা করা হয়। এখানে \texttt{color:red;} এবং \texttt{float:left;}--দুটি declaration সঠিকভাবে লেখা হয়েছে।
\end{enumerate}

\subsection{Practice MCQ}

\begin{enumerate}
\item HTML document-এর বিভিন্ন অংশ আলাদা করতে কোন tag ব্যবহৃত হয়?\\
ক) \texttt{<span>} \quad খ) \texttt{<div>} \quad গ) \texttt{<select>} \quad ঘ) \texttt{<option>}

\item Text-এর নির্দিষ্ট inline অংশ style করতে কোন tag বেশি উপযোগী?\\
ক) \texttt{<div>} \quad খ) \texttt{<span>} \quad গ) \texttt{<table>} \quad ঘ) \texttt{<hr>}

\item \texttt{<div>} কোন ধরনের element?\\
ক) Inline-level \quad খ) Block-level \quad গ) Empty tag \quad ঘ) Image tag

\item \texttt{<span>} কোন ধরনের element?\\
ক) Inline-level \quad খ) Block-level \quad গ) Table tag \quad ঘ) Empty tag

\item নিচের কোন style attribute syntax সঠিক?\\
ক) \texttt{<div style="color:red, float=left">}\\
খ) \texttt{<div style="color=red; float:left">}\\
গ) \texttt{<div style="color:red; float:left;">}\\
ঘ) \texttt{<div style="color:red" "float:left">}

\item \texttt{background-color} property-এর কাজ কী?\\
ক) text color পরিবর্তন করা\\
খ) background color পরিবর্তন করা\\
গ) link তৈরি করা\\
ঘ) image যোগ করা
\end{enumerate}

\begin{boardbox}{MCQ Answer Key}
১-খ, ২-খ, ৩-খ, ৪-ক, ৫-গ, ৬-খ
\end{boardbox}

\subsection{CQ Practice}

\begin{practicebox}
\textbf{উদ্দীপক:} রাফি তার কলেজের homepage তৈরি করছে। সে page-এর উপরে একটি header section, তার নিচে menu section এবং পরে main content section তৈরি করতে চায়। Main content-এর paragraph-এর ভিতরে ``HSC ICT students'' অংশটি red ও bold করতে চায়।

\textbf{প্রশ্ন}\\
ক. \texttt{<div>} tag কী?\\
খ. \texttt{<span>} tag inline element কেন?\\
গ. উদ্দীপকের জন্য HTML code লেখ।\\
ঘ. Webpage design-এ \texttt{<div>}, \texttt{<span>} ও \texttt{style} attribute-এর ভূমিকা বিশ্লেষণ কর।
\end{practicebox}

\subsubsection{গ. Code}

\begin{examplebox}{Code}
\begin{verbatim}
<!DOCTYPE html>
<html>
<head>
  <title>ABC College</title>
</head>
<body>
  <div style="background-color:navy; color:white; padding:15px; text-align:center;">
    <h1>ABC College</h1>
    <p>Official Website</p>
  </div>

  <div style="background-color:#eeeeee; padding:10px;">
    <a href="index.html">Home</a> |
    <a href="about.html">About</a> |
    <a href="notice.html">Notice</a> |
    <a href="contact.html">Contact</a>
  </div>

  <div style="padding:15px; border:1px solid gray; margin-top:10px;">
    <h2>Welcome</h2>
    <p>
      This website is designed for
      <span style="color:red; font-weight:bold;">HSC ICT students</span>.
    </p>
  </div>
</body>
</html>
\end{verbatim}
\end{examplebox}

\subsubsection{ঘ. বিশ্লেষণ}

\texttt{<div>} webpage-এর বড় section যেমন header, menu, content ও footer আলাদা করতে সাহায্য করে। \texttt{<span>} paragraph বা line-এর নির্দিষ্ট ছোট অংশ style করতে ব্যবহৃত হয়। \texttt{style} attribute দিয়ে সরাসরি CSS property-value ব্যবহার করে color, background, padding, margin, border, alignment ইত্যাদি নিয়ন্ত্রণ করা যায়। তাই এই তিনটি ব্যবহার করলে webpage বেশি সাজানো, আকর্ষণীয় ও ব্যবহারবান্ধব হয়।

\subsection{Common Mistake}

\begin{mistakebox}{সাধারণ ভুল}
\begin{itemize}
\item নির্দিষ্ট word style করতে \texttt{<div>} ব্যবহার করা।
\item পুরো section group করতে \texttt{<span>} ব্যবহার করা।
\item \texttt{style="color=red"} লেখা; সঠিক \texttt{style="color:red;"}।
\item CSS declaration-এর মাঝে comma ব্যবহার করা।
\item \texttt{background color} লেখা; সঠিক \texttt{background-color}।
\item semicolon \texttt{;} না দেওয়া।
\item \texttt{<div>} ও \texttt{<span>} closing tag না দেওয়া।
\item inline CSS বেশি ব্যবহার করে code অগোছালো করা।
\end{itemize}
\end{mistakebox}

\subsection{১ মিনিটে রিভিশন}

\begin{recapbox}
\texttt{<div>} বড় section বা block group করে। \texttt{<span>} text-এর ছোট inline অংশ group করে। \texttt{<div>} block-level এবং \texttt{<span>} inline-level। \texttt{style} attribute দিয়ে inline CSS লেখা যায়। Style syntax হলো \texttt{property:value;}। একাধিক style declaration semicolon দিয়ে আলাদা করতে হয়। Professional design-এ বারবার inline style না লিখে internal বা external CSS ব্যবহার করা ভালো।
\end{recapbox}


\begin{boardbox}{পরবর্তী টপিক}
পরের টপিক: Full Webpage Designing Practical--Header, Menu, Image, Table, List ও Link দিয়ে সম্পূর্ণ page তৈরি।
\end{boardbox}
